\documentclass[11pt]{article}
\usepackage{amsmath,amssymb,amsthm}
\usepackage[margin=1in]{geometry}
\newtheorem{theorem}{Theorem}
\newtheorem{lemma}{Lemma}
\newtheorem{definition}{Definition}
\title{V104: Canonical Affine-First Hybrid Compression}
\author{NC0\_k-Avoid Laboratory}
\date{}
\begin{document}
\maketitle

\begin{definition}[Canonical affine-first preprocessing]
For every output choose the minority target value, breaking ties toward zero.
Scan output blocks and retain a block iff its canonical affine-hull equations
increase the current rank. Let the resulting rank be $R$ and protect every
input variable occurring in a retained equation. Then scan the remaining
outputs and greedily select the first canonical functional anchor whose head is
unprotected, unused, and preserves acyclicity. Let $f$ be the number of selected
heads and define
\[
\eta_{AF}=n-R-f.
\]
\end{definition}

\begin{lemma}[Protected rank remains on roots]
Every retained affine equation is supported on the final root variables of the
functional DAG.
\end{lemma}
\begin{proof}
Every variable occurring in a retained affine equation is protected before the
functional phase, and protected variables are forbidden as functional heads.
\end{proof}

\begin{theorem}[Canonical affine-first avoider]
For $C:\{0,1\}^n\to\{0,1\}^m$ with $m>n$, the canonical preprocessing is
polynomial time and constructs a word outside $\operatorname{Range}(C)$ in
\[
O(2^{\eta_{AF}}\operatorname{poly}(N)).
\]
\end{theorem}
\begin{proof}
The functional phase leaves $n-f$ roots. By the lemma, the retained affine
system still has rank $R$ on those roots, so exactly
$2^{n-f-R}=2^{\eta_{AF}}$ root assignments survive. Each extends uniquely
through the total functional graph relations.

If $s_A$ affine output blocks were retained, then $s_A\le R$ because each
retained block increases rank. Hence
\[
m-s_A-f\ge m-R-f>n-R-f=\eta_{AF}.
\]
Evaluate the remaining original gates on every relaxed assignment and choose a
residual output word not observed. Restoring the selected canonical target bits
gives a missing original output because every exact selected fiber is contained
in its affine-hull or functional-graph relaxation. Empty canonical fibers or an
inconsistent canonical affine system give immediate missing target patterns.
\end{proof}

\begin{theorem}[Strict infinite-family separation]
For every $k\ge1$ there is a connected exact-stretch circuit with $n=8k$ and
$m=n+1$ for which the canonical algorithm itself obtains
\[
\eta_{AF}=3,
\]
while
\[
\lambda=8k,\qquad
\mu=k+3,\qquad
3k\le\beta\le7k,\qquad
\nu=4k+1.
\]
\end{theorem}
\begin{proof}
Use a cyclic $0x1e$ block on the first $4k$ variables and the V103 $0x16$
rank block on the second $4k$ variables, plus one internal and one cross-block
$0x17$ majority output. The affine-first phase retains exactly the $4k-1$
independent $0x16$ parity-hull blocks, so $R=4k-1$ and the second variable group
is protected. The $0x1e$ and $0x17$ canonical fibers have full affine hull.
The functional scan therefore accepts the first $4k-2$ cyclic $0x1e$ graph
relations and rejects the final two by acyclicity; majority has no functional
anchor. Thus $f=4k-2$ and
\[
\eta_{AF}=8k-(4k-1)-(4k-2)=3.
\]
The previous-parameter identities follow from the V97 no-peel support,
V101 maximum-head count, V102 strong-backdoor lower/upper bounds, and V103
canonical-hull rank calculation established for the same family.
\end{proof}

\paragraph{Boundary.}
The computable parameter $\eta_{AF}$ can still be linear in the worst case. No
unrestricted polynomial-time $NC^0_3$ avoidance algorithm, improved general
worst-case exponent, new circuit lower bound, novelty, peer review, or
P-versus-NP resolution is claimed.
\end{document}
